% !TEX root = task-3-3-4.tex
% Editable task bank. Append new BankTask blocks at the end of the document.
\documentclass[12pt,a4paper]{article}
\usepackage{physics-textbook}
\usepackage{graphicx,needspace}
\usetikzlibrary{decorations.pathmorphing}
\newif\ifShowSolutions
\ifdefined\TutorEdition
  \ShowSolutionstrue
\else
  \ShowSolutionsfalse % Change to \ShowSolutionstrue to display solutions.
\fi
\newcounter{banktask}
\newsavebox{\BankTaskBox}
\newcommand{\BankTaskContent}[3]{%
  \noindent\textbf{\thebanktask.}\quad #2\par
  \ifShowSolutions
    \Needspace{50mm}% Keep the solution heading with its opening text or diagram.
    #3\par
    \ifdefined\TutorEdition
      \ifcsname TutorSource#1\endcsname
        \nopagebreak\smallskip{\small\csname TutorSource#1\endcsname\par}%
      \fi
    \fi
  \fi
}
\newcommand{\BankTask}[3]{%
  \refstepcounter{banktask}\par\medskip
  \sbox{\BankTaskBox}{\begin{minipage}{\linewidth}%
    \BankTaskContent{#1}{#2}{#3}\end{minipage}}%
  \ifdim\dimexpr\ht\BankTaskBox+\dp\BankTaskBox\relax>\dimexpr\textheight-20mm\relax
    \Needspace{35mm}\BankTaskContent{#1}{#2}{#3}%
  \else
    \Needspace{\dimexpr\ht\BankTaskBox+\dp\BankTaskBox\relax}%
    \noindent\usebox{\BankTaskBox}\par
  \fi
  \medskip
}
\begin{document}
\BookSubsection{3.3.4}{Принцип суперпозиции электрических полей}


% Source 7D9e2D: https://ege.fipi.ru/bank/questions.php?search=1&proj=BA1F39653304A5B041B656915DC36B38&qid=7D9e2D
\BankTask{7D9e2D}{
Две маленькие закреплённые проводящие бусинки, расположенные в точках $A$ и $B$, несут на себе заряды $+3q>0$ и $-1{,}5q<0$ соответственно (см. рисунок). Точка $C$ находится посередине между бусинками $A$ и $B$.
\begin{enumerate}
\item На бусинку $B$ со стороны бусинки $A$ действует сила Кулона, направленная горизонтально вправо.
\item Напряжённость результирующего электростатического поля в точке $C$ направлена горизонтально влево.
\item Модули сил Кулона, действующих на бусинки, одинаковы.
\item Если бусинки соединить тонкой стальной проволокой, они начнут отталкиваться друг от друга.
\item Если бусинки соединить незаряженной стеклянной палочкой, их заряды станут одинаковыми.
\end{enumerate}
\begin{center}
\begin{tikzpicture}[book diagram]
\draw[densely dashed] (0,0)--(4,0);
\fill (0,0) circle (2pt) node[above=7pt] {$+3q$} node[below=5pt] {$A$};
\fill (2,0) circle (2pt) node[below=5pt] {$C$};
\fill (4,0) circle (2pt) node[above=7pt] {$-1{,}5q$} node[below=5pt] {$B$};
\end{tikzpicture}
\end{center}
}{
\smallskip\noindent\textbf{Решение от GPT}\par
1) Неверно: разноимённые заряды притягиваются, поэтому сила на бусинку $B$ направлена влево. 2) Неверно: в точке $C$ оба поля направлены вправо — от положительного заряда к отрицательному. 3) Верно по третьему закону Ньютона: $F_{AB}=F_{BA}$. 4) Верно: при соединении проводником суммарный положительный заряд $3q-1{,}5q=1{,}5q$ перераспределится, и бусинки будут отталкиваться. 5) Неверно: стекло — диэлектрик и не обеспечивает перетекания зарядов между бусинками. Ответ: $3,\ 4$.
}

\end{document}
